\documentclass[../../main.tex]{subfiles}
\begin{document}

\begin{flushright}
  \footnotesize\textit{【自動文字起こし・要確認】}
\end{flushright}

\noindent
{\bf [解答]} 題意から
\begin{align}
  \begin{cases}
    2\beta = \alpha + \gamma \quad \dots \text{①} \\
    \sin^2\beta = \sin\alpha \sin\gamma \quad \dots \text{②}
  \end{cases}
\end{align}
①を②に代入
\begin{align}
  \sin^2 \dfrac{\alpha+\gamma}{2} &= 2\sin\dfrac{\alpha}{2}\cos\dfrac{\alpha}{2} \sin\dfrac{\gamma}{2}\cos\dfrac{\gamma}{2} \\
  &= \dfrac{1}{2} \left[ \sin\dfrac{\alpha+\gamma}{2} + \sin\dfrac{\alpha-\gamma}{2} \right] \left[ \sin\dfrac{\alpha+\gamma}{2} - \sin\dfrac{\alpha-\gamma}{2} \right]
\end{align}
$t = \dfrac{\alpha+\gamma}{2}, s = \dfrac{\alpha-\gamma}{2}$ とおく．
\begin{align}
  \sin^2 t + \sin^2 s = 0 \quad \dots \text{③}
\end{align}
$\sin t, \sin s \in \mathbb{R}$ から，③が成立するのは
\begin{align}
  \sin t = \sin s = 0
\end{align}
\begin{align}
  \therefore \dfrac{\alpha+\gamma}{2} = n\pi, \quad \dfrac{\alpha-\gamma}{2} = k\pi \quad (n, k \in \mathbb{Z})
\end{align}
の時である．解いて
\begin{align}
  \alpha = (n+k)\pi, \quad \gamma = (n-k)\pi
\end{align}
①に代入して
\begin{align}
  \beta = n\pi
\end{align}
従って，題意のようになるのは，$n, k \in \mathbb{Z}$ として
\begin{align}
  (\alpha, \beta, \gamma) = ((n+k)\pi, n\pi, (n-k)\pi)
\end{align}

\newpage

\end{document}
